% Mathématiques 6e — cours V3
% Initiation à la pensée informatique

\section{Objectifs}

\begin{itemize}
\item Identifier une instruction et une séquence d'instructions.
\item Produire puis exécuter une séquence d'instructions.
\item Repérer les données d'entrée et le résultat de sortie d'une procédure simple.
\item Répéter manuellement une même séquence pour accomplir une tâche.
\item Reconnaître qu'une répétition peut remplacer plusieurs copies identiques d'une séquence.
\item Programmer la construction d'un trajet simple sur une grille.
\item Exécuter un algorithme pas à pas pour vérifier son résultat.
\item Repérer et corriger une erreur simple dans une séquence.
\end{itemize}

\section{1. Penser avant de programmer}

La pensée informatique ne commence pas par un langage de programmation.

Elle commence par une question :
\[
\text{« Quelles instructions faut-il exécuter, et dans quel ordre ? »}
\]

Une procédure peut être décrite oralement, sur papier, avec des flèches, par des blocs ou dans un programme.

\begin{remarque}
En 6e, le cœur du travail est la compréhension des instructions, des séquences et des répétitions. Il n'est pas nécessaire de transformer immédiatement chaque activité en code textuel complexe.
\end{remarque}

\section{2. Une instruction}

Une instruction est une action suffisamment précise pour pouvoir être exécutée.

Par exemple :
\begin{itemize}
\item avancer d'une case ;
\item tourner à droite ;
\item tracer un segment ;
\item afficher un message ;
\item prendre la valeur donnée en entrée.
\end{itemize}

\begin{definition}
Une \textbf{instruction} décrit une action élémentaire que l'exécutant peut réaliser sans avoir à deviner ce qui est demandé.
\end{definition}

L'exécutant peut être un élève, un robot, un logiciel ou un personnage dans un environnement de programmation par blocs.

\section{3. Une séquence d'instructions}

Une séquence est une suite ordonnée d'instructions.

L'ordre est essentiel.

\begin{exemple}
Considérons les instructions :
\begin{enumerate}
\item avancer de deux cases ;
\item tourner à droite ;
\item avancer d'une case.
\end{enumerate}

Exécuter :
\[
1,\quad2,\quad3
\]
ne produit pas nécessairement le même trajet que :
\[
2,\quad1,\quad3.
\]
\end{exemple}

Modifier l'ordre peut modifier le résultat.

\coursefigure{/images/mathematiques/6eme/informatique-sequence.svg}{Robot sur une grille et séquence d'instructions représentée par des flèches.}{Une séquence doit être exécutée dans l'ordre prévu ; chaque instruction transforme la position ou l'orientation obtenue à l'étape précédente.}

\section{4. Entrée et sortie}

Une procédure peut recevoir une information au départ.

On parle d'\textbf{entrée}.

Elle produit ensuite un résultat.

On parle de \textbf{sortie}.

\begin{exemple}
Une procédure reçoit en entrée un nombre.

Elle demande :
\begin{enumerate}
\item ajouter $4$ ;
\item doubler le résultat.
\end{enumerate}

Avec l'entrée :
\[
3,
\]
on obtient :
\[
3+4=7,
\]
puis :
\[
2\times7=14.
\]

La sortie vaut :
\[
\boxed{14}.
\]
\end{exemple}

\section{5. Exécuter pas à pas}

Pour comprendre un algorithme, on peut suivre l'état après chaque instruction.

\begin{tabular}{c|c}
Étape & Valeur \\
entrée & $3$ \\
après l'ajout & $7$ \\
après le double & $14$ \\
\end{tabular}

Cette méthode évite de perdre l'information intermédiaire.

\begin{methode}
Pour exécuter une séquence :
\begin{enumerate}
\item noter l'état initial ;
\item lire une seule instruction ;
\item l'exécuter ;
\item noter le nouvel état ;
\item passer à l'instruction suivante ;
\item continuer jusqu'à la fin.
\end{enumerate}
\end{methode}

\section{6. Décrire un trajet sur une grille}

Une grille permet de représenter un déplacement de manière simple.

Un robot peut posséder :
\begin{itemize}
\item une position ;
\item une orientation ;
\item des instructions de déplacement.
\end{itemize}

Par exemple :
\begin{verbatim}
avancer
avancer
tourner à droite
avancer
\end{verbatim}

Chaque instruction est interprétée à partir de la position et de l'orientation actuelles.

\section{7. Position et orientation}

Deux robots situés sur la même case peuvent réagir différemment à l'instruction :
\[
\text{« avancer »}
\]
s'ils ne regardent pas dans la même direction.

L'état du robot ne se limite donc pas toujours à sa position.

On doit parfois connaître :
\[
\text{position}+\text{orientation}.
\]

\begin{attention}
Une séquence de déplacement peut sembler correcte sur le papier mais échouer si l'orientation initiale n'est pas précisée.
\end{attention}

\section{8. Construire une séquence à partir d'un objectif}

On peut partir du résultat attendu et chercher les instructions nécessaires.

\begin{exemple}
Un robot part en bas à gauche d'une grille et regarde vers la droite.

La cible se trouve trois cases à droite puis deux cases vers le haut.

Une séquence possible est :
\begin{verbatim}
avancer
avancer
avancer
tourner à gauche
avancer
avancer
\end{verbatim}
\end{exemple}

La séquence doit être testée pour vérifier qu'elle atteint réellement la cible.

\section{9. Une même tâche peut avoir plusieurs séquences}

Deux algorithmes peuvent parfois atteindre le même objectif avec des instructions différentes.

Une solution peut être :
\begin{itemize}
\item plus courte ;
\item plus facile à comprendre ;
\item plus facile à vérifier ;
\item plus simple à modifier.
\end{itemize}

L'objectif en 6e n'est pas toujours de trouver la solution la plus courte, mais de produire une séquence correcte, compréhensible et testable.

\section{10. Répéter une séquence}

Certaines tâches utilisent plusieurs fois la même suite d'actions.

Supposons que l'on doive répéter quatre fois :
\begin{verbatim}
avancer
tourner à droite
\end{verbatim}

On peut exécuter manuellement la même séquence :
\[
4
\]
fois.

Cette répétition permet par exemple de suivre les quatre côtés d'un carré si les longueurs et les angles sont adaptés.

\section{11. Comprendre une répétition}

\coursefigure{/images/mathematiques/6eme/informatique-repetition.svg}{Comparaison entre quatre copies d'une séquence et une structure de répétition unique.}{Une répétition décrit qu'une même séquence doit être exécutée plusieurs fois ; elle réduit les recopies et met en évidence la structure du problème.}

On peut symboliquement écrire :
\begin{verbatim}
répéter 4 fois
    avancer
    tourner à droite
\end{verbatim}

La structure « répéter » ne change pas les actions réalisées : elle les organise plus efficacement.

\section{12. Répétition manuelle}

En 6e, on doit être capable d'exécuter une répétition sans ordinateur.

\begin{exemple}
La séquence :
\begin{verbatim}
avancer de 2 cases
tourner à droite
\end{verbatim}
doit être répétée :
\[
3
\]
fois.

On déroule :
\begin{enumerate}
\item avancer de deux cases, tourner ;
\item avancer de deux cases, tourner ;
\item avancer de deux cases, tourner.
\end{enumerate}

La répétition a donc exécuté :
\[
3
\]
déplacements et :
\[
3
\]
rotations.
\end{exemple}

\section{13. Compter le nombre d'exécutions}

Une erreur fréquente consiste à répéter une fois de trop ou une fois de moins.

Si une séquence est exécutée :
\[
5
\]
fois et qu'elle contient :
\[
2
\]
instructions, le nombre total d'instructions élémentaires exécutées vaut :
\[
5\times2=10,
\]
si les deux instructions sont exécutées à chaque répétition.

Cette relation permet de contrôler le déroulement.

\section{14. Construire un trajet simple}

Programmer un trajet simple consiste à transformer une route dessinée ou imaginée en une séquence précise.

\begin{methode}
\begin{enumerate}
\item repérer le départ ;
\item repérer l'orientation initiale ;
\item localiser la cible ;
\item découper le trajet en segments simples ;
\item écrire les déplacements et rotations ;
\item exécuter la séquence ;
\item corriger si la cible n'est pas atteinte.
\end{enumerate}
\end{methode}

\section{15. Exemple développé : contourner un obstacle}

Un robot se déplace sur une grille.

Un obstacle bloque la case située juste devant lui.

Il peut utiliser :
\begin{verbatim}
tourner à droite
avancer
tourner à gauche
avancer
tourner à gauche
avancer
tourner à droite
\end{verbatim}

La bonne séquence dépend de la forme de l'obstacle et de l'orientation initiale.

L'algorithme doit être testé sur la grille réelle.

\section{16. Décomposer un problème}

Une tâche complexe devient plus simple lorsqu'on la décompose.

Par exemple :
\[
\text{aller de A à D}
\]
peut être découpé en :
\[
A\longrightarrow B,
\]
puis :
\[
B\longrightarrow C,
\]
puis :
\[
C\longrightarrow D.
\]

Chaque sous-trajet peut être résolu séparément avant de réunir les séquences.

\begin{propriete}
Décomposer une tâche en sous-tâches permet de réduire la difficulté et facilite la vérification de chaque étape.
\end{propriete}

\section{17. Prévoir avant d'exécuter}

Avant de lancer une séquence sur une machine, il est utile de prévoir le résultat.

Cette prévision permet ensuite de comparer :
\[
\text{résultat attendu}
\]
et
\[
\text{résultat obtenu}.
\]

Si les deux diffèrent, il faut chercher l'origine de l'écart.

\section{18. Déboguer une séquence}

Déboguer signifie chercher et corriger une erreur dans une procédure ou un programme.

\begin{methode}
\begin{enumerate}
\item prévoir ce qui devrait se passer ;
\item exécuter la séquence pas à pas ;
\item repérer la première étape où le résultat devient incorrect ;
\item identifier l'instruction responsable ;
\item modifier cette instruction ;
\item recommencer le test depuis le début.
\end{enumerate}
\end{methode}

Cette démarche est plus efficace que modifier plusieurs instructions au hasard.

\section{19. Exemple développé : une instruction incorrecte}

Un robot doit aller deux cases à droite puis une case vers le haut.

Il regarde initialement vers la droite.

La séquence proposée est :
\begin{verbatim}
avancer
avancer
tourner à droite
avancer
\end{verbatim}

Après deux déplacements, le robot est correctement placé horizontalement.

Mais tourner à droite l'oriente vers le bas.

La première erreur apparaît donc à la troisième instruction.

Il faut remplacer :
\[
\text{« tourner à droite »}
\]
par :
\[
\text{« tourner à gauche »}.
\]

\section{20. Lire un programme par blocs}

Dans un environnement comme Scratch, une instruction peut être représentée par un bloc.

Plusieurs blocs emboîtés forment une séquence.

Un bloc « répéter » contient les instructions à exécuter plusieurs fois.

Le principe mathématique reste le même que sur papier :
\[
\text{instruction}
\longrightarrow
\text{séquence}
\longrightarrow
\text{répétition}.
\]

L'interface graphique change, mais le raisonnement reste transférable.

\section{21. Entrées et sorties dans un programme}

Une entrée peut être :
\begin{itemize}
\item une position de départ ;
\item une orientation ;
\item un nombre choisi ;
\item une longueur ;
\item une consigne donnée par l'utilisateur.
\end{itemize}

Une sortie peut être :
\begin{itemize}
\item une position finale ;
\item un dessin ;
\item un nombre obtenu ;
\item un message affiché.
\end{itemize}

Identifier entrée et sortie aide à comprendre ce que réalise réellement la procédure.

\section{22. Variables et conditions : seulement comme prolongement}

Certains environnements de programmation utilisent des variables ou des conditions.

Ces notions peuvent apparaître lors d'une activité, mais elles ne constituent pas le cœur des objectifs de ce chapitre de 6e.

\begin{remarque}
Le socle attendu ici reste : identifier et produire des instructions et séquences, répéter une séquence et programmer un trajet simple.
\end{remarque}

Cette distinction évite de transformer l'initiation de 6e en cours formel de programmation.

\section{23. Exemple développé : construire un carré}

Un robot graphique doit tracer un carré.

Une procédure possible consiste à répéter quatre fois :
\begin{verbatim}
avancer de la même longueur
tourner d'un angle droit
\end{verbatim}

Un angle droit mesure :
\[
90^\circ.
\]

La répétition comporte :
\[
4
\]
déplacements identiques et :
\[
4
\]
rotations de :
\[
90^\circ.
\]

Au total, la rotation accumulée vaut :
\[
4\times90^\circ=360^\circ.
\]

Le robot retrouve donc son orientation initiale.

\section{24. Exemple développé : prédire une sortie numérique}

La procédure suivante reçoit un nombre :
\begin{verbatim}
ajouter 5
doubler
\end{verbatim}

Avec l'entrée :
\[
7,
\]
on obtient :
\[
7+5=12,
\]
puis :
\[
2\times12=24.
\]

La sortie est :
\[
\boxed{24}.
\]

Pour vérifier la compréhension, on peut ensuite essayer une autre entrée sans changer la séquence.

\section{25. Méthode générale}

\begin{methode}
Pour résoudre un problème de pensée informatique :
\begin{enumerate}
\item identifier l'objectif final ;
\item préciser les données d'entrée ;
\item décomposer la tâche ;
\item écrire les instructions dans l'ordre ;
\item repérer les séquences répétées ;
\item exécuter ou simuler pas à pas ;
\item comparer le résultat obtenu au résultat attendu ;
\item corriger la première erreur détectée ;
\item tester à nouveau.
\end{enumerate}
\end{methode}

\section{26. Erreurs fréquentes}

\begin{itemize}
\item Donner une consigne trop vague pour être exécutée.
\item Oublier que l'ordre des instructions compte.
\item Ne pas préciser l'orientation initiale d'un robot.
\item Répéter une séquence une fois de trop.
\item Confondre le nombre de répétitions et le nombre total d'instructions exécutées.
\item Modifier plusieurs étapes simultanément pendant un débogage.
\item Regarder uniquement la position finale sans rechercher la première instruction erronée.
\item Présenter variables et conditions comme des prérequis alors qu'elles ne sont pas au cœur des attendus de 6e.
\end{itemize}

\section{À retenir}

\begin{aretenir}
Une instruction décrit une action précise.

Une séquence est une suite ordonnée d'instructions.

Une procédure peut recevoir des entrées et produire une sortie.

Lorsqu'une même séquence doit être exécutée plusieurs fois, on peut la répéter manuellement ou utiliser une structure de répétition.

En 6e, il faut notamment savoir :
\begin{itemize}
\item identifier une instruction ou une séquence ;
\item produire et exécuter une séquence ;
\item répéter manuellement une séquence pour accomplir une tâche ;
\item programmer la construction d'un trajet simple.
\end{itemize}

Le débogage consiste à exécuter pas à pas afin de repérer la première étape qui ne correspond plus au résultat attendu.
\end{aretenir}
